%==================================================================
% SAMPUL LUAR  (Bahasa Indonesia — mengikuti contoh hal. 1)
% Datanya dari a0-identitas.tex
%==================================================================
\thispagestyle{empty}
\begin{center}
    \ifdefstring{\jenisnaskah}{proposal}{%
        % ---- PROPOSAL: label "PROPOSAL TUGAS AKHIR" di ATAS judul (mengikuti contoh) ----
        \MakeUppercase{\labeltipe}\\[1.5cm]
        \begin{spacing}{1.5}
            \textbf{\large \MakeUppercase{\judulid}}\\[1.5cm]
        \end{spacing}
    }{%
        % ---- TUGAS AKHIR: judul di ATAS label ----
        \begin{spacing}{1.5}
            \textbf{\large \MakeUppercase{\judulid}}\\[1.5cm]
        \end{spacing}
        \textbf{\MakeUppercase{\labeltipe}}\\[0.3cm]
        \ifdefempty{\submitline}{}{\textbf{\submitline}\\[1.5cm]}
    }

    % Logo diatur dengan HEIGHT (bukan width) agar tinggi sampul tetap
    % sama untuk logo berbentuk bujur sangkar maupun memanjang.
    \includegraphics[height=3.2cm]{logo-institusi}\\[1.5cm]

    \textbf{Disusun oleh:}\\
    \textbf{\MakeUppercase{\penulis}}\\
    \textbf{\nim}

    \vfill

    \textbf{\MakeUppercase{Program Studi \prodi}}\\
    \textbf{\MakeUppercase{\fakultas}}\\
    \textbf{\MakeUppercase{\universitas}}\\
    \textbf{\MakeUppercase{JAKARTA}}\\
    \textbf{\tahun}
\end{center}
\clearpage
